The orchestration rule in Boreal is simple: match the lead first, plan the work second, and decompose only when needed. This rule is intentionally conservative. It rejects the assumption that the first useful response to an ambiguous request is always a task tree. In many requests, especially higher-complexity ones, the decisive question is not ``what are the steps'' but ``who or what should lead the work.''

Lead-first orchestration has two benefits. The first is quality. Once a credible lead \texttt{Supply} is identified, the system can shape the rest of the plan around a realistic execution path. The second is safety. Premature decomposition creates pressure to allocate sub-work before approval, before commercial boundaries are clear, and before the system knows whether the fulfillment path is digital, embodied, or hybrid.

This matters directly for embodied work. Suppose a request involves an onsite visit and a follow-up report. A tool-first planner may start by drafting a checklist, an email template, and a report outline because those are callable outputs. A lead-first planner should instead ask whether the request requires local presence, whether a field-capable supply exists in the right geography, and whether the owner has specified the access and verification conditions. Only after that lead route is credible should the system derive sub-work such as documentation, summary writing, or operator review.

The same logic applies to hybrid work. A local human may need to perform inspection, pickup, installation, or witness tasks, while an agent or tool handles background research, document assembly, or structured reconciliation. Lead-first orchestration allows the digital steps to support the embodied path rather than silently replacing it. In this sense, the rule is not only a planning heuristic. It is a guardrail against tool-shaped underplanning.
